\section{Данные и подготовка выборки}

\subsection{Источник данных}

В работе используются данные \textbf{USDA NASS CSB/CDL} за 2017--2024 годы. CDL предоставляет годовые классы культур, а
CSB-слой используется как источник полевых идентификаторов, площади и пространственного контекста. Такой набор данных
позволяет перейти от анализа отдельных наблюдений к анализу историй использования конкретных полей.

Выбор данных USDA NASS CSB/CDL обусловлен их открытостью, масштабом и воспроизводимостью. Данный набор предоставляет
последовательные сведения о культурах на уровне полей за несколько лет и тем самым позволяет проверить методологию
построения рекомендательной системы на устойчивой экспериментальной базе. Для российской территории открытый
стандартизованный аналог уровня «поле --- год --- культура» отсутствует в публичном доступе. Поэтому данные USDA NASS
CSB/CDL рассматриваются как экспериментальная основа для разработки и проверки подхода. Перенос предложенной
методологии на российский контекст требует локальных данных, адаптации агрономических правил и дополнительной
экспертной валидации.

На уровне исходной записи данные содержат:

\begin{itemize}
\tightlist
\item
  идентификатор поля;
\item
  площадь поля;
\item
  региональные признаки;
\item
  пространственные координаты;
\item
  годовые значения культурных классов за несколько последовательных лет.
\end{itemize}

\subsection{Предобработка исходных данных}

На первом этапе исходные данные были приведены к виду, пригодному для последующего моделирования. Основная задача
предобработки состояла в том, чтобы перейти от исходных CDL-кодов к более интерпретируемому представлению культур и
удалить классы, не соответствующие целевой постановке исследования. Для каждого года выполнялось сопоставление кодов с
названиями культур и их укрупненными группами. После этого из выборки исключались нерелевантные для задачи классы,
например неаграрные или неопределённые категории. Итогом этого этапа стал очищенный датасет, сохранённый как
промежуточный артефакт воспроизводимого конвейера.

По метаданным артефакта после базовой очистки выборка имеет следующие характеристики:

\begin{itemize}
\tightlist
\item
  8 110 870 строк;
\item
  33 столбца;
\item
  8 годовых CDL-столбцов, соответствующих годам 2017--2024.
\end{itemize}

Предобработка на этом этапе была необходима не только для очистки исходного массива, но и для формирования
единообразного представления культур, пригодного для дальнейшего анализа временных последовательностей.

\subsection{Формирование оконной выборки}

После очистки исходная таблица была преобразована в оконный датасет, который используется для обучения модели. Вместо
того чтобы работать с отдельной строкой как с изолированным наблюдением, каждая запись была преобразована в окно вида

\begin{equation}
\notag
(c_{t-3}, c_{t-2}, c_{t-1}) \rightarrow c_t
\end{equation}

где \(c_{t-3}\), \(c_{t-2}\) и \(c_{t-1}\) соответствуют столбцам \texttt{history\_1}, \texttt{history\_2}, \texttt{history\_3}, а \(c_t\) --- целевой культуре \texttt{target}.

Пример нескольких строк такого представления приведён в таблице~\ref{tab:window-dataset-example}. В ней видно, что
одна строка соответствует одному полю и одному прогнозируемому году: технический идентификатор и
пространственно-региональные
признаки сохраняют контекст поля, три предыдущих значения культуры используются как входная история, а значение в
столбце \texttt{target} является следующей культурой, которую должна предсказать модель.

\begin{table}[H]
\centering
\small
\caption{Пример формирования оконной выборки}
\label{tab:window-dataset-example}
\begin{tabular}{rllllr}
\toprule
\texttt{CSBID} &
\texttt{history\_1} &
\texttt{history\_2} &
\texttt{history\_3} &
\texttt{target} &
\texttt{target\_year} \\
\midrule
\texttt{351724000000036} & \texttt{wheat} & \texttt{wheat} & \texttt{sorghum} & \texttt{sorghum} & 2020 \\
\texttt{351724000000037} & \texttt{wheat} & \texttt{wheat} & \texttt{sorghum} & \texttt{sorghum} & 2020 \\
\texttt{351724000000044} & \texttt{forage\_hay} & \texttt{forage\_hay} & \texttt{other\_cereals} & \texttt{fallow} & 2020 \\
\bottomrule
\end{tabular}
\end{table}

Поскольку в исходных данных доступны годы с 2017 по 2024, из одной исходной
записи можно сформировать несколько перекрывающихся окон. В результате полный оконный датасет содержит 40 554 350 строк.

Преимущество такой схемы состоит в том, что она позволяет использовать историю поля как явный вход модели. Тем самым
задача переходит от статического анализа признаков к анализу краткой временной последовательности культур.

\subsection{Фильтрация редких целевых классов}

После построения оконного датасета был выполнен анализ распределения целевой переменной и фильтрация редких
целевых классов. Порог редкости был выбран на уровне 1\% от общего числа наблюдений. Такой порог позволяет сохранить
основные массовые классы и одновременно исключить категории, для которых число
наблюдений недостаточно для устойчивой оценки качества.

После фильтрации:

\begin{itemize}
\tightlist
\item
  в выборке осталось 38 511 210 строк;
\item
  итоговое пространство целевых классов было сокращено до 9 культурных классов.
\end{itemize}

Итоговый набор целевых классов имеет вид:

\begin{table}[htbp]
\centering
\caption{Итоговые и исключённые целевые классы}
\label{tab:target-classes}
\begin{tabular}{ll}
\toprule
Оставшиеся целевые классы & Исключённые редкие классы \\
\midrule
\texttt{corn} & \texttt{oilseeds\_other} \\
\texttt{cotton} & \texttt{sugar\_crops} \\
\texttt{fallow} & \texttt{peanuts} \\
\texttt{forage\_hay} & \texttt{rice} \\
\texttt{legumes} & \texttt{root\_tubers} \\
\texttt{other\_cereals} & \texttt{sunflower} \\
\texttt{sorghum} & \texttt{specialty\_crops} \\
\texttt{soybeans} & \texttt{vegetables\_melons} \\
\texttt{wheat} & \texttt{fruits\_berries} \\
\bottomrule
\end{tabular}
\end{table}

Агрегация и последующая фильтрация целевых классов являются важной частью постановки. С одной стороны, они уменьшают
разреженность и дисбаланс данных. С другой стороны, они ограничивают область применимости системы: модель работает не с
полным пространством исходных культур, а с устойчивым и интерпретируемым набором укрупнённых целевых классов.

\subsection{Разделение на обучающую, валидационную и тестовую выборки}

После формирования итоговой выборки было выполнено единое разделение на обучающую, валидационную и тестовую выборки.
Это разделение фиксируется один раз и затем переиспользуется во всех основных этапах работы.
Такой подход необходим для
воспроизводимости результатов: все модели и все сравнительные эксперименты должны оцениваться на одном и том же
разбиении.

Принципиально важно, что разбиение выполняется с учётом группировки по \texttt{CSBID}: окна, построенные
для одного и того же поля, не должны одновременно попадать в разные части выборки. Такой контракт предотвращает утечку
информации между разными частями выборки и делает оценку качества более корректной.

Фактические размеры выборок зафиксированы в таблице~\ref{tab:data-preparation-sizes}.

\begin{center}
\small
\begin{longtable}{p{0.27\textwidth} r p{0.20\textwidth} p{0.27\textwidth}}
\caption{Размеры выборок на этапах подготовки данных}
\label{tab:data-preparation-sizes}\\
\toprule
Этап & Размер, строк & Источник & Комментарий \\
\midrule
\endfirsthead

\toprule
Этап & Размер, строк & Источник & Комментарий \\
\midrule
\endhead

\bottomrule
\endlastfoot

Очищенный набор данных &
8 110 870 &
\texttt{01\_meta.json} &
Таблица после первичной подготовки и очистки данных \\

Оконная выборка &
40 554 350 &
\texttt{02\_meta.json} &
Выборка формата «история культур \(\rightarrow\) целевая культура» \\

Выборка после фильтрации &
38 511 210 &
\texttt{02\_meta.json} &
Оконная выборка после исключения редких целевых классов \\

Обучающая выборка &
26 957 630 &
\texttt{run\_meta.json} &
Часть фиксированного разбиения с учётом группировки по \texttt{CSBID} \\

Валидационная выборка &
5 776 566 &
\texttt{run\_meta.json} &
Часть фиксированного разбиения для подбора и промежуточной оценки \\

Тестовая выборка &
5 777 014 &
\texttt{run\_meta.json} &
Часть фиксированного разбиения для итоговой оценки \\

\end{longtable}
\end{center}

Таким образом, на данном этапе формируется единая воспроизводимая база для всех последующих экспериментов: 
обучения базовых и основных моделей, оценки рекомендаций, анализа уверенности и интерпретации результатов.

\subsection{Итоговое признаковое представление}

На основе подготовленной оконной выборки было сформировано итоговое признаковое представление,
используемое на последующих этапах моделирования.
Оно сочетает историю культур, региональный контекст и базовые пространственные характеристики поля.

Состав признаков и их типы приведены в таблице~\ref{tab:final-model-features}.

\begin{center}
\small
\begin{longtable}{p{0.24\textwidth} p{0.23\textwidth} p{0.43\textwidth}}
\caption{Итоговое признаковое представление}
\label{tab:final-model-features}\\
\toprule
Признак & Тип & Краткий смысл \\
\midrule
\endfirsthead

\toprule
Признак & Тип & Краткий смысл \\
\midrule
\endhead

\bottomrule
\endlastfoot

\texttt{history\_1} &
Категориальный &
Культура на поле в году \(t-3\) \\

\texttt{history\_2} &
Категориальный &
Культура на поле в году \(t-2\) \\

\texttt{history\_3} &
Категориальный &
Культура на поле в году \(t-1\), ближайшем к целевому году \\

\texttt{STATEFIPS} &
Категориальный &
Код штата, задающий региональный контекст наблюдения \\

\texttt{CNTYFIPS} &
Категориальный &
Код округа, задающий более локальный региональный контекст \\

\texttt{CSBACRES} &
Числовой &
Площадь поля по данным CSB \\

\texttt{INSIDE\_X} &
Числовой &
Пространственная координата поля по оси X \\

\texttt{INSIDE\_Y} &
Числовой &
Пространственная координата поля по оси Y \\

\end{longtable}
\end{center}

В результате подготовки данных была сформирована единая оконная выборка, связывающая историю культур на поле,
пространственно-региональный контекст и целевую культуру следующего года. Полученное признаковое представление далее
используется для обучения предиктивных моделей, формирования рекомендательного списка и анализа интерпретируемости
результатов.
